\begin{ex}
Trong bảng các số từ 1 đến 100, số đứng ngay sau số 99 là số nào?
\shortans{100}
\loigiai{
Số đứng liền sau số 99 là số $100$.
\textbf{Đáp số:} $100$.
}
\end{ex}

\begin{ex}
Trong bảng các số từ 1 đến 100, có bao nhiêu số có hai chữ số?
\shortans{90}
\loigiai{
Các số có hai chữ số từ 10 đến 99 có: $99 - 10 + 1 = 90$ (số).
\textbf{Đáp số:} $90$.
}
\end{ex}

\begin{ex}
Trong bảng các số từ 1 đến 100, có bao nhiêu số có một chữ số?
\shortans{9}
\loigiai{
Từ 1 đến 100 có các số một chữ số là: 1, 2, 3, 4, 5, 6, 7, 8, 9 (tất cả 9 số).
\textbf{Đáp số:} $9$.
}
\end{ex}

\begin{ex}
Số bé nhất có ba chữ số trong bảng là số nào?
\shortans{100}
\loigiai{
Số 100 là số duy nhất có ba chữ số trong bảng các số từ 1 đến 100.
\textbf{Đáp số:} $100$.
}
\end{ex}

\begin{ex}
Trong cùng một cột của bảng số từ 1 đến 100, hai số ở hai hàng liền nhau hơn kém nhau bao nhiêu đơn vị?
\shortans{10}
\loigiai{
Trong bảng số, mỗi hàng có 10 số, do đó hai ô ở hai hàng liền nhau trên cùng một cột hơn kém nhau đúng $10$ đơn vị.
\textbf{Đáp số:} $10$.
}
\end{ex}

\begin{ex}
Số ở ô ngay phía dưới số 45 trong bảng số là số nào?
\shortans{55}
\loigiai{
Số ở ô ngay dưới cùng cột cách số ở trên 10 đơn vị: $45 + 10 = 55$.
\textbf{Đáp số:} $55$.
}
\end{ex}

\begin{ex}
Số ở ô ngay phía trên số 68 trong bảng số là số nào?
\shortans{58}
\loigiai{
Số ở ô ngay trên cùng cột kém số ở dưới 10 đơn vị: $68 - 10 = 58$.
\textbf{Đáp số:} $58$.
}
\end{ex}

\begin{ex}
Số ở ô ngay bên phải của số 37 trong bảng số là số nào?
\shortans{38}
\loigiai{
Số ở ngay bên phải là số liền sau: $37 + 1 = 38$.
\textbf{Đáp số:} $38$.
}
\end{ex}

\begin{ex}
Số ở ô ngay bên trái của số 80 trong bảng số là số nào?
\shortans{79}
\loigiai{
Số ở ngay bên trái là số liền trước: $80 - 1 = 79$.
\textbf{Đáp số:} $79$.
}
\end{ex}

\begin{ex}
Hàng thứ 5 trong bảng các số từ 1 đến 100 gồm các số từ 41 đến bao nhiêu?
\shortans{50}
\loigiai{
Hàng thứ 5 gồm 10 số: $41, 42, 43, 44, 45, 46, 47, 48, 49, 50$.
\textbf{Đáp số:} $50$.
}
\end{ex}

\begin{ex}
Cột cuối cùng bên phải của bảng số gồm toàn các số tròn chục. Số cuối cùng ở góc dưới bên phải là số nào?
\shortans{100}
\loigiai{
Số ở góc dưới cùng bên phải của bảng số là số $100$.
\textbf{Đáp số:} $100$.
}
\end{ex}

\begin{ex}
Có bao nhiêu số có chữ số tận cùng là 7 trong bảng các số từ 1 đến 100?
\shortans{10}
\loigiai{
Các số đó là: $7, 17, 27, 37, 47, 57, 67, 77, 87, 97$. Có tất cả $10$ số.
\textbf{Đáp số:} $10$.
}
\end{ex}

\begin{ex}
Có bao nhiêu số có chữ số hàng chục là 5 trong bảng các số từ 1 đến 100?
\shortans{10}
\loigiai{
Các số đó là hàng từ 50 đến 59: $50, 51, \dots, 59$. Có tất cả $10$ số.
\textbf{Đáp số:} $10$.
}
\end{ex}

\begin{ex}
Tìm số đứng liền trước của số 100.
\shortans{99}
\loigiai{
Số liền trước của 100 là: $100 - 1 = 99$.
\textbf{Đáp số:} $99$.
}
\end{ex}

\begin{ex}
Tìm số tự nhiên $x$ biết $x$ đứng giữa số 88 và số 90 trong bảng số.
\shortans{89}
\loigiai{
Số đứng giữa 88 và 90 là số $89$.
\textbf{Đáp số:} $89$.
}
\end{ex}

\begin{ex}
Đếm cách 5 đơn vị: $5, 10, 15, 20, \dots, 35$. Số thích hợp điền vào chỗ chấm là số nào?
\shortans{25}
\loigiai{
Đếm thêm 5 đơn vị: $20 + 5 = 25$.
\textbf{Đáp số:} $25$.
}
\end{ex}

\begin{ex}
Đếm cách 10 đơn vị: $10, 20, 30, \dots, 60$. Số thích hợp điền vào chỗ chấm là số nào?
\shortans{40}
\loigiai{
Đếm thêm 10 đơn vị: $30 + 10 = 40$.
\textbf{Đáp số:} $40$.
}
\end{ex}

\begin{ex}
Trong bảng số, từ số 21 đến số 30 có tất cả bao nhiêu số?
\shortans{10}
\loigiai{
Số lượng số từ 21 đến 30 là: $30 - 21 + 1 = 10$ (số).
\textbf{Đáp số:} $10$.
}
\end{ex}

\begin{ex}
Trong bảng số, từ số 1 đến số 50 có bao nhiêu số tròn chục?
\shortans{5}
\loigiai{
Các số tròn chục là: $10, 20, 30, 40, 50$. Có tất cả $5$ số.
\textbf{Đáp số:} $5$.
}
\end{ex}

\begin{ex}
Tìm một số biết rằng số đó ở ngay bên dưới số 72 và ở ngay bên trái số 83.
\shortans{82}
\loigiai{
Số ở dưới số 72 là: $72 + 10 = 82$. Số ở bên trái số 83 là: $83 - 1 = 82$. Vậy số đó là $82$.
\textbf{Đáp số:} $82$.
}
\end{ex}

\begin{ex}
Một hàng có 10 ô số. Bạn Hùng đã điền được 6 ô số. Hỏi Hùng còn phải điền bao nhiêu ô số nữa để hoàn thành hàng đó?
\shortans{4}
\loigiai{
Số ô số Hùng còn phải điền là: $10 - 6 = 4$ (ô số).
\textbf{Đáp số:} $4$.
}
\end{ex}

\begin{ex}
Một bảng số bị rách một góc làm mất các số từ 95 đến 100. Hỏi có bao nhiêu số bị mất?
\shortans{6}
\loigiai{
Số các số bị mất là: $100 - 95 + 1 = 6$ (số).
\textbf{Đáp số:} $6$.
}
\end{ex}

\begin{ex}
Số lớn nhất có hai chữ số khác nhau trong bảng số là số nào?
\shortans{98}
\loigiai{
Số lớn nhất có 2 chữ số khác nhau là $98$.
\textbf{Đáp số:} $98$.
}
\end{ex}

\begin{ex}
Tổng của chữ số hàng chục và hàng đơn vị của số 96 bằng bao nhiêu?
\shortans{15}
\loigiai{
Tổng hai chữ số là: $9 + 6 = 15$.
\textbf{Đáp số:} $15$.
}
\end{ex}

\begin{ex}
Hiệu giữa chữ số hàng chục và hàng đơn vị của số 73 bằng bao nhiêu?
\shortans{4}
\loigiai{
Hiệu hai chữ số là: $7 - 3 = 4$.
\textbf{Đáp số:} $4$.
}
\end{ex}
